\begin{table*}[htbp]
\centering
\small
\caption{Устойчивость решения к десяти преобразованиям текста по четырём процедурам}
\label{tab:7}
\begin{tabular}{>{\raggedright\arraybackslash}p{0.131\linewidth}>{\raggedright\arraybackslash}p{0.131\linewidth}>{\raggedright\arraybackslash}p{0.131\linewidth}>{\raggedright\arraybackslash}p{0.131\linewidth}>{\raggedright\arraybackslash}p{0.131\linewidth}>{\raggedright\arraybackslash}p{0.131\linewidth}>{\raggedright\arraybackslash}p{0.131\linewidth}}
\toprule
Преобразование & Вход признаков не изменился & Процедура 2: нестабильность & Процедура 2: смена решения & Процедура 1: доля $|$Δ$|$ $>$ 5 пунктов & Процедура 3: медиана $|$ΔNLL$|$ & Процедура 4: доля $|$Δ$|$ $>$ 5 пунктов \\
\midrule
t11 расширение текста & 0.000 & 0.633 & 0.451 & 0.567 & 0.4366 & 0.233 \\
t04 перестановка абзацев & 0.250 & 0.089 & 0.055 & 0.267 & 0.0279 & 0.133 \\
t02 снятие markdown & 0.533 & 0.126 & 0.023 & 0.067 & 0.0000 & 0.033 \\
t05 исправление пунктуации & 0.017 & 0.055 & 0.017 & 0.100 & 0.0066 & 0.050 \\
t10 сокращение текста & 0.000 & 0.269 & 0.015 & 0.567 & 0.1584 & 0.117 \\
t03 срезание заголовков & 0.600 & 0.051 & 0.003 & 0.000 & 0.0000 & 0.000 \\
t13 опечатки & 0.000 & 0.051 & 0.001 & 0.050 & 0.1207 & 0.000 \\
t15 гомоглифы & 0.000 & 0.002 & 0.001 & 0.000 & 0.0216 & 0.033 \\
t01 замена тире & 0.083 & 0.000 & 0.000 & 0.000 & 0.0019 & 0.033 \\
t16 перенос между форматами & 0.450 & 0.143 & 0.000 & 0.217 & 0.0000 & 0.050 \\
\bottomrule
\end{tabular}
\begin{minipage}{\linewidth}\footnotesize\vspace{2pt}
\textit{Примечание.} Единица анализа — у процедуры 2 — пара «документ × допустимая модель»; у процедур 1, 3 и 4 — документ. Знаменатель: 60 документов панели; 464 пары на преобразование у процедуры 2, 60 документов на преобразование у остальных. Данные: процедура 2 — стресс-ревизия r5, процедура 4 — r5, процедуры 1 и 3 — r4 в области пригодности десяти преобразований. Доли и медианы по ячейкам; интервалы не считались. Событие нестабильности определяется правилом своей процедуры, поэтому колонки между собой не складываются. Шкала: меньше значит устойчивее решение. Сокращения: Δ — сдвиг оценки после преобразования; NLL — negative log-likelihood; смена решения — переход через порог 0.5; «вход признаков не изменился» — доля документов, у которых после преобразования совпали профиль prose, профиль full и все счётчики разметки, то есть процедура получила прежний вход. Пропуски: преобразование t14 исключено: амендмент ревизии r5 перевёл его в статус not executable, поэтому знаменатель равен десяти, а не одиннадцати. Поправка на множественные сравнения не применялась. Стресс-тест описательный.
\end{minipage}
\end{table*}
